We define the graph stream and its access constraints, specify the condensed
memory and its allocation, and then describe how memory utility is evaluated.

\subsection{Graph Stream and Historical Access}
\label{sec:stream_setting}

Citation networks and Web product catalogs receive new documents or products
in successive batches. Storage costs and data retention policies can limit
access to earlier records for model updates. We model this situation by
allowing construction access to each raw batch for a limited period, during
which useful training information must be retained in a smaller memory.

The stream has $T$ stages. The \emph{current increment} at stage $t$ is
$\Delta_t=(V_t,X_t,T_t,Y_t,E_t)$. Here $V_t$ contains the $n_t=|V_t|$
\emph{construction nodes}, $X_t\in\mathbb R^{n_t\times d_x}$ contains their
$d_x$-dimensional features, $T_t=\{t_i:i\in V_t\}$ contains their texts,
and $Y_t=\{y_i:i\in V_t\}$ contains their available labels. The edges
$E_t\subseteq V_t\times V_t$ are induced on these nodes, with adjacency
$A_t\in\{0,1\}^{n_t\times n_t}$. Neighborhood information used in
construction comes from the current increment. Held-out nodes are excluded from this input so their
features and texts cannot serve as construction examples.

Let $\mathcal C_t$ denote the classes present in the current increment, and let
$\mathcal C_{\leq t}=\bigcup_{\tau=1}^t\mathcal C_\tau$ denote all
classes observed so far. Our primary evaluation uses class-incremental streams,
following continual graph condensation studies~\cite{cat,puma}. Their class sets
are disjoint, and a \emph{class order} determines which classes enter each stage.
This partition tests retention as the label space expands. We also evaluate
chronological streams, in which publication year determines arrival and
classes may recur across stages. \autoref{sec:setup} gives both schedules.

At stage $t$, the constructor receives $\Delta_t$ and the
\emph{historical memory} $M_{t-1}$. Once the new memory is saved, the raw
nodes, texts, features, and edges of $\Delta_t$ are unavailable to later
construction stages. Memory therefore carries earlier information into later
construction. The goal is to retain training information for predicting all
observed classes using this compact memory. Downstream graph access is
specified separately below.

\subsection{Condensed Memory and Fixed Condensation Rate}
\label{sec:memory_setting}

A \emph{memory item} contains a numerical feature, readable text, and a
class label. A \emph{memory slot} is the position allocated to an item
before its feature and text are constructed. The new block is
$M_t^+=\{(x_j^+,t_j^+,y_j)\}_{j=1}^{m_t}$, where $m_t$ is the number of
new slots and superscript $+$ denotes the new block. Starting from
$M_0=\varnothing$, the constructor $\mathcal F_t$ at stage $t$ produces
\begin{equation}
 M_t=M_{t-1}\oplus M_t^+,
 \qquad M_t^+=\mathcal F_t(\Delta_t,M_{t-1}),
\label{eq:continual_memory}
\end{equation}
where $\oplus$ appends the new items. Historical memory remains fixed
during this operation. This design preserves previously constructed features
and texts when their sources are unavailable, and separates construction of
new items from historical revision. We use this fixed access policy
throughout the experiments.

The \emph{fixed condensation rate} $\rho\in(0,1]$ controls the number of retained
items relative to the number of observed construction nodes. With
$N_t=\sum_{\tau=1}^{t}n_\tau$, the nominal capacity is
$\bar K_t=\operatorname{round}(\rho N_t)$. We reserve at least one new
slot for each class in $\mathcal C_t$, giving actual capacity
$K_t=\max\{\bar K_t,|M_{t-1}|+|\mathcal C_t|\}$ and
$m_t=K_t-|M_{t-1}|$. Remaining slots use largest-remainder rounding,
subject to the number of available nodes in each class. This prevents
integer rounding from excluding a class in the current increment. We report the
effective rate $K_t/N_t$ when this coverage requirement changes the nominal
rate. Following graph condensation with a fixed rate~\cite{cat}, the allocation
keeps compression comparable as data accumulates. Memory
size grows with the stream. The rate constrains the number of retained
nodes rather than bytes or GPU memory. Its numerical value and
sensitivity study are given with the experiments.

For \method{}, the \emph{support matrix}
$Q\in\mathbb R_{\geq0}^{n_t\times m_t}$ links current nodes to new slots.
Its learnable scores are $Z\in\mathbb R^{n_t\times m_t}$, and
$R\in\mathbb R^{m_t\times d_x}$ denotes the feature residual added to
the pooled features. Text $\bar t_i$ denotes the cleaned source string,
and $e_i\in\mathbb R^{d_e}$ its normalized embedding, where $d_e$ is the
text encoder's output dimension. \emph{Hard text selection} chooses a
source string for a slot. A \emph{soft text representation} is a weighted embedding mixture
used to supply gradients. \emph{Materialization} means saving the final
features, selected strings, labels, and source identifiers. The main
memory is \emph{edge-free}, storing feature and text items without edges
between them while incorporating graph context during construction.
\autoref{sec:method_interface} explains the
corresponding computations.

\subsection{Downstream Learning and Evaluation}
\label{sec:evaluation_protocol}

The \emph{student model} is the downstream predictor trained on condensed
memory using graph and text pairing and class prompt supervision. A \emph{class prompt} is a
predefined text description of a label. Prediction at stage $t$ compares a
node's graph representation
$h_i^g$ with the projected, normalized prompt representation $p_c$ for
$c\in\mathcal C_{\leq t}$, selecting the largest cosine similarity. This
readout with class prompts evaluates whether memory supports graph and language
alignment~\cite{graphclip}. Construction and training use class labels,
so the task is supervised classification of observed classes. A
\emph{proxy model} supplies a training signal during memory construction.

Evaluation is transductive. The student model can use the original graph
structure and node features at inference, while evaluation labels are
withheld from construction and training. An \emph{evaluation mask}
specifies which nodes are scored. Matched comparisons use common text
processing, encoders, class prompts, memory allocation, class orders,
student models, and downstream training budgets. The implementation,
metrics, and repetition statistics are specified in~\autoref{sec:setup}.
